% Academic Research Paper — Standard Two-Column Format
\documentclass[10pt,twocolumn]{article}

\usepackage[margin=0.75in]{geometry}
\usepackage{amsmath,amssymb}
\usepackage{booktabs}
\usepackage{cite}
\usepackage[dvipsnames]{xcolor}
\usepackage{hyperref}

\hypersetup{colorlinks=true,linkcolor=MidnightBlue,citecolor=MidnightBlue,urlcolor=MidnightBlue}

\begin{document}

% ZONE:TITLE_AUTHORS:START
\title{\textbf{Scalable Multi-Agent Document Generation and Synthesis}}
\author{
  \textbf{First Author}$^{1}$, \textbf{Second Author}$^{2}$ \\
  $^{1}$Department of Computer Science, University of Technology \\
  $^{2}$AI Research Laboratories, Institute of Science \\
  \texttt{\{first.author, second.author\}@institution.edu}
}
\date{\today}
\maketitle
% ZONE:TITLE_AUTHORS:END

% ZONE:ABSTRACT:START
\begin{abstract}
Automated document synthesis requires balancing structural consistency, domain-specific formatting constraints, and responsive multi-turn human-in-the-loop editing. In this paper, we present an end-to-end multi-agent pipeline for modular LaTeX compilation and repair. Our architecture introduces isolated zone agents responsible for targeted document sections, coordinated by a stateful orchestrator that ensures compilation invariants. Empirical evaluation on benchmark corpora demonstrates a 42\% reduction in compilation faults and a 3.2x speedup in document revision latency compared to monolithic approaches.
\end{abstract}
% ZONE:ABSTRACT:END

% ZONE:INTRODUCTION:START
\section{Introduction}
Document authoring in technical domains frequently relies on typesetting systems such as \LaTeX{} to produce publication-ready manuscripts. However, integrating Large Language Models (LLMs) into iterative document editing introduces major hurdles: broken control sequences, syntax hallucinations, and full-document re-generation costs.

In this work, we propose a modular architecture that solves these challenges:
\begin{itemize}
  \item \textbf{Zone Decomposition}: Parsing documents into distinct, independently editable functional zones with strict boundary isolation.
  \item \textbf{Specialist Agents}: Routing localized edits to dedicated domain experts that enforce semantic and structural schemas.
  \item \textbf{Resilient Compilation}: Auto-repairing engine that intercepts compiler errors and dynamically injects missing packages, color definitions, and font mappings.
\end{itemize}
% ZONE:INTRODUCTION:END

% ZONE:RELATED_WORK:START
\section{Related Work}
Prior investigations into automated document generation primarily focus on monolithic prompting techniques~\cite{lamport1994latex,knuth1984texbook}. While effective for simple snippets, monolithic generation suffers from cumulative hallucination rates on multi-page manuscripts.

Recent advances in multi-agent workflows have demonstrated improved performance across coding and mathematical reasoning tasks~\cite{vaswani2017attention}. We extend these paradigms to structured typesetting and incremental AST maintenance.
% ZONE:RELATED_WORK:END

% ZONE:METHODOLOGY:START
\section{Methodology}
Our system consists of three core components: the Zone Engine, the Specialist Agent Pool, and the Compatibility Engine.

\subsection{Zone Representation}
Given a document $\mathcal{D}$, we partition it into an ordered set of zones $\mathcal{Z} = \{z_1, z_2, \dots, z_n\}$ such that:
\begin{equation}
\mathcal{D} = \mathcal{H} \cup \left( \bigcup_{i=1}^n z_i \right) \cup \mathcal{F}
\end{equation}
where $\mathcal{H}$ represents the preamble and document start, and $\mathcal{F}$ represents the document terminus.
% ZONE:METHODOLOGY:END

% ZONE:EXPERIMENTS:START
\section{Experiments}
We benchmark our approach across 1,000 synthetic editing prompts and evaluate compilation success rate and latency.

\begin{table}[ht]
\centering
\caption{Performance comparison on document revision benchmarks.}
\label{tab:results}
\begin{tabular}{lrr}
\toprule
\textbf{Approach} & \textbf{Success Rate (\%)} & \textbf{Latency (s)} \\
\midrule
Monolithic LLM   & 74.2 & 14.8 \\
Section Patcher  & 88.5 & 6.2 \\
Our System       & \textbf{98.9} & \textbf{2.1} \\
\bottomrule
\end{tabular}
\end{table}
% ZONE:EXPERIMENTS:END

% ZONE:DISCUSSION:START
\section{Discussion}
The quantitative improvements highlighted in Table~\ref{tab:results} stem directly from zone isolation. By restricting the LLM context to the targeted zone, we minimize hallucinated package collisions and preserve unchanged sections verbatim.
% ZONE:DISCUSSION:END

% ZONE:CONCLUSION:START
\section{Conclusion}
We presented a modular multi-agent framework for resilient \LaTeX{} document synthesis. Our architecture enables fast, reliable document authoring with automated error repair. Future work includes expanding real-time collaborative editing and multimodal figure synthesis.
% ZONE:CONCLUSION:END

% ZONE:REFERENCES:START
\begin{thebibliography}{99}
\bibitem{lamport1994latex}
L.~Lamport, \emph{\LaTeX: A Document Preparation System}, 2nd~ed.\hskip 1em plus 0.5em minus 0.4em\relax Addison-Wesley, 1994.

\bibitem{knuth1984texbook}
D.~E.~Knuth, \emph{The \TeX{}book}.\hskip 1em plus 0.5em minus 0.4em\relax Addison-Wesley, 1984.

\bibitem{vaswani2017attention}
A.~Vaswani \emph{et al.}, ``Attention is all you need,'' in \emph{Advances in Neural Information Processing Systems (NeurIPS)}, 2017, pp. 5998--6008.
\end{thebibliography}
% ZONE:REFERENCES:END

\end{document}
